% ============================================================
%  Formalizing Quantitative Equational Reasoning
%  Program Verification Exam -- 07/09/2026
%
%  30-minute version.  Each definition is followed immediately by
%  its Rocq encoding.  The long deck (with the
%  [frel_derives_weakening] walkthrough, the metric-view system,
%  and the enriched-Lawvere framing) is in presentation-long.tex.
%
%  Compile with:  latexmk -pdf presentation.tex
% ============================================================
\documentclass[10pt,aspectratio=169]{beamer}

\usetheme[progressbar=frametitle,block=fill]{metropolis}

\usepackage[T1]{fontenc}
\usepackage[utf8]{inputenc}
\usepackage{amsmath,amssymb,amsfonts}
\usepackage{bm}
\usepackage{stmaryrd}          % \llbracket \rrbracket
\usepackage{mathtools}
\usepackage{mathpartir}        % inference rules
\usepackage{listings}
\usepackage{xcolor}
\usepackage[most]{tcolorbox}
\usepackage{tikz}
\usetikzlibrary{arrows.meta,positioning,decorations.pathmorphing}
\usepackage{booktabs}
\usepackage{nicematrix}           % row/column-labelled context matrices
\usepackage{appendixnumberbeamer}  % \inserttotalframenumber excludes
% the appendix

% ============================================================
%  Listings setup for Rocq source and for Rocq goal states
% ============================================================

\definecolor{rocqkw}{HTML}{5B4A7A}     % declaration heads only
\definecolor{rocqcon}{HTML}{B3541E}    % constructors / our names
\definecolor{rocqcom}{HTML}{6A737D}    % comments
\definecolor{goalbg}{HTML}{F4F6F8}
\definecolor{goalrule}{HTML}{C4CCD4}
\definecolor{srcbg}{HTML}{FBF9F4}

\lstdefinelanguage{Rocq}{
  % Only the declaration heads are marked; everything else stays plain.
  morekeywords=[1]{Record, Inductive, Definition, Fixpoint, Lemma, Theorem,
    Proposition, Corollary, Notation, Section, End, Variable, Hypothesis,
    Context, Proof, Qed, Defined},
  sensitive=true,
  morecomment=[n]{(*}{*)},
  morestring=[b]",
}

\lstdefinestyle{rocqsrc}{
  language=Rocq,
  basicstyle=\ttfamily\fontsize{6.6pt}{7.5pt}\selectfont,
  keywordstyle=[1]\color{rocqkw},
  commentstyle=\color{rocqcom}\itshape,
  showstringspaces=false,
  columns=fullflexible,
  keepspaces=true,
  breaklines=true,
  breakatwhitespace=false,
  literate={->}{{$\rightarrow$}}2 {=>}{{$\Rightarrow$}}2
           {<=}{{$\leq$}}2 {>=}{{$\geq$}}2 {|-}{{$\vdash$}}2
           {\\/}{{$\vee$}}2 {/\\}{{$\wedge$}}2,
  aboveskip=0.2em, belowskip=0.2em,
}

\lstdefinestyle{rocqgoal}{
  style=rocqsrc,
  basicstyle=\ttfamily\fontsize{6pt}{6.9pt}\selectfont,
}

\lstnewenvironment{rocq}[1][]{\lstset{style=rocqsrc,#1}}{}
\lstnewenvironment{rocqg}[1][]{\lstset{style=rocqgoal,#1}}{}

% A framed box for a proof-state dump
% A proof state, split so the conclusion is unmistakable:
%   upper part = hypotheses (context), lower part = what must be proved.
\definecolor{goalconcbg}{HTML}{E3EDF6}
\definecolor{goalconcfg}{HTML}{0C3149}

\lstnewenvironment{gctx}[1][]{\lstset{style=rocqgoal,#1}}{}
\lstnewenvironment{gconc}[1][]{%
  \lstset{style=rocqgoal,
    basicstyle=\ttfamily\fontsize{6.4pt}{7.3pt}\selectfont\color{goalconcfg},#1}}{}

\newtcolorbox{goalbox}[1][]{
  bicolor,
  colback=goalbg, colbacklower=goalconcbg,
  colframe=goalrule,
  boxrule=0.4pt, arc=1pt,
  left=2pt, right=2pt, top=1.5pt, bottom=1.5pt,
  middle=1.5pt,
  before skip=3pt, after skip=3pt,
  fonttitle=\scriptsize\bfseries,
  #1
}

\newtcblisting{srcbox}[1][]{
  listing only,
  listing options={style=rocqsrc},
  colback=srcbg, colframe=goalrule,
  boxrule=0.4pt, arc=1pt,
  left=2pt, right=2pt, top=1.5pt, bottom=1.5pt,
  before skip=3pt, after skip=3pt,
  fonttitle=\scriptsize\bfseries,
  #1
}

% ============================================================
%  Stepped proof scripts: a script pane whose "current" line is
%  highlighted, earlier lines normal, later lines dimmed.
%
%    \pstep{3}                 -- declare which step this slide shows
%    \pline{0}{...}            -- context line, never dimmed
%    \pline{n}{...}            -- line belonging to step n
% ============================================================

\definecolor{stepbg}{HTML}{FBEDC6}     % the line you are on
\definecolor{stepdim}{HTML}{A9B1B9}    % not executed yet

\newcommand{\pcur}{0}
\newcommand{\pstep}[1]{\renewcommand{\pcur}{#1}}
% a strut exactly one \baselineskip tall (so it tracks the pane's font
% size), making consecutive highlighted lines merge into one block
\newcommand{\pstrut}{\vrule height 0.76\baselineskip depth 0.24\baselineskip width 0pt}
\newcommand{\pline}[2]{%
  \ifnum#1=0\relax
    \makebox[\linewidth][l]{\pstrut#2}%
  \else\ifnum#1<\pcur\relax
    \makebox[\linewidth][l]{\pstrut#2}%
  \else\ifnum#1=\pcur\relax
    {\setlength{\fboxsep}{0pt}%
     \colorbox{stepbg}{\makebox[\linewidth][l]{\pstrut#2}}}%
  \else
    \makebox[\linewidth][l]{\pstrut\color{stepdim}#2}%
  \fi\fi\fi\par}

\newtcolorbox{scriptbox}[1][]{
  colback=srcbg, colframe=goalrule,
  boxrule=0.4pt, arc=1pt,
  left=2pt, right=2pt, top=1.5pt, bottom=1.5pt,
  before skip=3pt, after skip=3pt,
  fontupper=\ttfamily\fontsize{6pt}{7.6pt}\selectfont\frenchspacing,
  #1
}


\newcommand{\ptag}[1]{\hfill\textit{#1}}

% ------------------------------------------------------------
%  The proof script of [kantorovich_models_ica] (Proposition 4.6).
%    step 1 -- setup, case split into the four axioms
%    step 2 -- the three equational axioms
%  Then the interpolative axiom, one tactic per step:
%    step 3 -- evaluate the two [+_p] terms
%    step 4 -- convexity of the Kantorovich lifting
%    step 5 -- monotonicity of the weighted sum
%    step 6 -- non-expansiveness of the interpretation
%    step 7 -- read the (0,2) entry off the context matrix
%    step 8 -- the (1,3) component, symmetric
% ------------------------------------------------------------
\newcommand{\modelscript}{%
    \pline{0}{Theorem kantorovich\_models\_ica :}
    \pline{0}{\ \ fuzzy\_models kantorovich\_algebra (@ica\_theory R).}
    \pline{1}{Proof.}
    \pline{1}{\ \ rewrite /fuzzy\_models.}
    \pline{1}{\ \ move => Y phi Hphi.}
    \pline{1}{\ \ case: Hphi => [p | p | p q | p eps delta ...] rho.}
    \pline{2}{\ \ - by rewrite kantorovich\_eval\_op}
    \pline{2}{\ \ \ \ \ \ convex\_mixture\_idem.\ptag{Idem}}
    \pline{2}{\ \ - by rewrite !kantorovich\_eval\_op}
    \pline{2}{\ \ \ \ \ \ convex\_mixture\_skew\_comm.\ptag{SkewComm}}
    \pline{2}{\ \ - by rewrite !kantorovich\_eval\_op}
    \pline{2}{\ \ \ \ \ \ convex\_mixture\_skew\_assoc.\ptag{SkewAssoc}}
    \pline{3}{\ \ - rewrite !kantorovich\_eval\_op.\ptag{Interp}}
    \pline{4}{\ \ \ \ apply: (le\_trans}
    \pline{4}{\ \ \ \ \ \ (kantorovich\_convex\_mixture \_ \_ \_ \_ \_ Hd)).}
    \pline{5}{\ \ \ \ apply: weighted\_sum\_le.}
    \pline{6}{\ \ \ \ - apply: (le\_trans (interpretation\_nexp}
    \pline{6}{\ \ \ \ \ \ \ \ rho (inord 0) (inord 2))).}
    \pline{7}{\ \ \ \ \ \ by rewrite /= /ica\_interp\_rel !inordK.}
    \pline{8}{\ \ \ \ - apply: (le\_trans (interpretation\_nexp}
    \pline{8}{\ \ \ \ \ \ \ \ rho (inord 1) (inord 3))).}
    \pline{8}{\ \ \ \ \ \ by rewrite /= /ica\_interp\_rel !inordK.}
    \pline{8}{Qed.}
}

% ------------------------------------------------------------
%  The proof script of [derives_fin_sound].
%    step 1 -- unfold, introduce the interpretation
%    step 2 -- the [remember] trick on the derivation mode
%    step 3 -- induction: twelve goals, one per rule
%    step 4 -- the Birkhoff cases
%    step 5 -- [FD_UseVariables], walked through on its own slide
%    step 6 -- the remaining cases
% ------------------------------------------------------------
\newcommand{\soundnessscript}{%
    \pline{0}{Theorem derives\_fin\_sound (A : FuzzyAlgebra R sig)}
    \pline{0}{\ \ \ \ (E : fuzzy\_theory R sig) (X : fuzzy\_space R)}
    \pline{0}{\ \ \ \ (phi : @judgement R sig (fcarrier X)) :}
    \pline{0}{\ \ fuzzy\_models A E -> derives\_fin E phi ->}
    \pline{0}{\ \ satisfies A X phi.}
    \pline{0}{Proof.}
    \pline{1}{\ \ rewrite /fuzzy\_models /satisfies /derives\_fin.}
    \pline{1}{\ \ move => Hmodel Hderive rho.}
    \pline{2}{\ \ remember FRelFinite as mode eqn:Hmode in Hderive.}
    \pline{2}{\ \ revert Hmode.}
    \pline{3}{\ \ induction Hderive; intros Hmode.}
    \pline{4}{\ \ - apply: Hmodel X phi H rho.\ptag{Init}}
    \pline{4}{\ \ - by [].\ptag{EqRefl}}
    \pline{4}{\ \ - symmetry. exact: IHHderive ...\ptag{EqSym}}
    \pline{4}{\ \ - by rewrite (IHHderive1 ...) ...\ptag{EqTrans}}
    \pline{4}{\ \ - congr fa\_ops; apply ...\ptag{EqCong}}
    \pline{5}{\ \ - move => //=.\ptag{UseVars}}
    \pline{5}{\ \ \ \ apply: (interpretation\_nexp rho).}
    \pline{6}{\ \ - apply: satisfies\_subst. ...\ptag{Subst}}
    \pline{6}{\ \ - rewrite (IHHderive1 ...). ...\ptag{QEqRepL}}
    \pline{6}{\ \ - rewrite -(IHHderive2 ...). ...\ptag{QEqRepR}}
    \pline{6}{\ \ - apply: le\_trans ... H.\ptag{Up}}
    \pline{6}{\ \ - apply: (andP (frel\_range ...)).2\ptag{Max}}
    \pline{6}{\ \ - discriminate Hmode.\ptag{OrderCompl}}
    \pline{6}{Qed.}
}

% ------------------------------------------------------------
%  Notation
% ------------------------------------------------------------
\NiceMatrixOptions{code-for-first-row=\scriptstyle,code-for-first-col=\scriptstyle}
% \ctx{<col spec>}{<body: first row = column labels, first col = row labels>}
\newcommand{\ctx}[2]{%
    \left(\;
        \begin{bNiceArray}{#1}[first-row,first-col,small]#2
\end{bNiceArray}\right)}
\newcommand{\qeq}[1]{=_{#1}}
\newcommand{\den}[1]{\llbracket #1 \rrbracket}
\newcommand{\ICA}{\mathbb{I}\mathbb{C}\mathbb{A}}
\newcommand{\Dist}{\mathcal{D}}
\newcommand{\Terms}{\mathbb{T}}
\newcommand{\eps}{\varepsilon}
\newcommand{\Kan}{\mathcal{K}}
\newcommand{\Cpl}{\Gamma}
\newcommand{\bpl}[1]{+_{#1}}
\newcommand{\file}[1]{{\color{rocqcon}\texttt{#1}}}
\newcommand{\rid}[1]{{\ttfamily\small #1}}

% metropolis: make frametitles and the title page a bit tighter
\metroset{sectionpage=none,numbering=fraction}
\makeatletter
\setlength{\metropolis@frametitle@padding}{1.1ex}
\makeatother
\setbeamerfont{title}{size=\Large,series=\bfseries}
\setbeamerfont{author}{size=\normalsize}
\setbeamerfont{institute}{size=\small}
\setbeamerfont{date}{size=\small}
\setbeamerfont{frametitle}{size=\normalsize,series=\bfseries}
\setbeamersize{text margin left=5mm,text margin right=5mm}

\title{Formalizing Quantitative Equational Reasoning}
\author{Benjamin Bennetzen}
\institute{Aalborg University \quad\textbullet\quad \texttt{bcb@aau.dk}\\[0.3em]
{\small Supervised by Giorgio Bacci and Robin Kaarsgaard}}
\date{Program Verification Exam \quad\textbullet\quad 07/09/2026}

\begin{document}

\begin{frame}[plain]
    \titlepage
\end{frame}

% ============================================================
\section{Setting the scene}
% ============================================================

\begin{frame}{Motivation}
    \begin{itemize}
        \item[--] Equations are fundamental to mathematical reasoning.
            \medskip
        \item[--] Reasoning about programs is also based on program
            equivalences.
            \medskip
        \item[--] But exact equivalence is often too strict. An
            approximation, errors, probabilistic and quantum computing.
            \medskip
        \item[--] We want quantitative analogues of equational
            reasoning.
    \end{itemize}
\end{frame}

\begin{frame}{Replace standard equations with quantitative ones}
    \begin{center}
        \begin{tcolorbox}[colback=srcbg,colframe=goalrule,width=0.55\textwidth,
            halign=center,boxrule=0.6pt,arc=8pt]
            \Huge $s \qeq{\eps} t$
        \end{tcolorbox}
    \end{center}
    \pause
    \vfill
    \begin{center}
        \large ``$s$ is within $\eps$ distance of $t$''
    \end{center}
    % \vfill
    % {\footnotesize To make this precise we need a syntax of terms, a
    %     semantics in which distances live, and a proof system relating the
    %     two. Today: that story for one running example, with the Rocq
    % encoding shown next to each definition.}
\end{frame}

% ============================================================
\section{The classical picture: barycentric algebras}
% ============================================================

\begin{frame}{An equational theory of barycentric algebras}
    \begin{columns}[T]
        \begin{column}{0.42\textwidth}
            \textbf{Variables}\\[0.3em]
            $s,t,u,\dots$
        \end{column}
        \begin{column}{0.55\textwidth}
            \textbf{Operations}\\[0.3em]
            $\Sigma_{CA} = \{\,\bpl{p} \mid 0 < p < 1\,\}$,\\
        \end{column}
    \end{columns}
    \vfill
    \textbf{Equations}
    \begin{align*}
        \text{(Idempotency)}\qquad & s \bpl{p} s = s\\[0.3em]
        \text{(Skew commutativity)}\qquad & s \bpl{p} t = t \bpl{1-p} s\\[0.3em]
        \text{(Skew associativity)}\qquad & (s \bpl{p} t) \bpl{q} u
        = s \bpl{pq} \bigl(t \bpl{\frac{(1-p)q}{1-pq}} u\bigr)
    \end{align*}
    \vfill
    {\footnotesize Also known as convex algebras or convex spaces.}
\end{frame}

\begin{frame}[fragile]{Formalization: signatures and terms}
    \file{src/Signature.v}
    \begin{srcbox}
Record signature := {
  sym : Type;
  arity : sym -> nat;
}.

Inductive term (sig : signature) (X : Type) : Type :=
  | Var : X -> term sig X
  | App : forall f : sym sig, ('I_(arity f) -> term sig X) -> term sig X.

Definition algebra_ops (sig : signature) (A : Type) : Type :=
  forall f : sym sig, (ordinal (arity f) -> A) -> A.
    \end{srcbox}
    \begin{itemize}
        \item \rid{sym} is an arbitrary \rid{Type}, so the signature may be
            uncountable, as $\Sigma_{\mathrm{CA}}$ is.
        \item Arguments are given as functions with argumnets of type ordinal
            \rid{'I\_n}.
    \end{itemize}
\end{frame}

\begin{frame}{A model of the theory}
    A model interprets each $\bpl{p}$ while respecting the axioms.
    \vfill
    One model: the finitely supported probability distributions
    $\Dist(X)$ with convex combinations
    \[
        (\mu \bpl{p} \nu)(x) \;=\; p\,\mu(x) + (1-p)\,\nu(x).
    \]
    \vfill
    \pause
    \textbf{Small example.} $X = \{a,b,c\}$, written as vectors
    $(\mu(a),\mu(b),\mu(c))$. Skew associativity with $p = q = \tfrac12$
    (so $pq = \tfrac14$ and $\tfrac{(1-p)q}{1-pq} = \tfrac13$):
    \[
        \underbrace{(\delta_a \bpl{\frac12} \delta_b) \bpl{\frac12} \delta_c}
        _{(\frac14,\frac14,\frac12)}
        \;=\;
        \underbrace{\delta_a \bpl{\frac14}
        (\delta_b \bpl{\frac13} \delta_c)}_{(\frac14,\frac14,\frac12)}
    \]
\end{frame}

\begin{frame}[fragile]{Formalization: distributions and convex mixtures}
    \begin{srcbox}
Record probability_weight (R : realType) := {
  weight : R;   weight_unit : (0 <= weight <= 1)%R;
}.
Record probability_distribution (R : realType) (X : finType) := {
  probability_mass :> X -> R;
  probability_mass_ge0 : forall x, (0 <= probability_mass x)%R;
  probability_mass_total : \sum_(x : X) probability_mass x = 1;
}.
Definition weighted_sum (p : probability_weight R) (x y : R) : R :=
  weight p * x + (1 - weight p) * y.

Definition convex_mixture (p : probability_weight R)
    (mu nu : probability_distribution R X) : probability_distribution R X.
Proof.
  refine {| probability_mass := fun x => weighted_sum p (mu x) (nu x) |}.
  ...
Defined.
    \end{srcbox}
    \begin{itemize}
        \item \rid{X : finType} makes ``finitely supported'' automatic.
        \item The three barycentric axioms become
            \rid{convex\_mixture\_idem}, \rid{\_skew\_comm} and
            \rid{\_skew\_assoc}: each is a scalar identity such as
            $px + (1-p)x = x$, discharged by \rid{lra} or \rid{field} and
            then lifted by \rid{probability\_distribution\_ext}.
    \end{itemize}
\end{frame}

% ============================================================
\section{Going quantitative}
% ============================================================

\begin{frame}{How to extend to the quantitative setting}
    An algebraic theory has three equivalent presentations.
    \vfill
    \begin{description}
        \item[\underline{Quantitative} equational theories]\hfill\\
            operations $+$ \underline{quantitative} equations.
            \medskip
        \item[\underline{EMet}-Lawvere theories]\hfill\\
            Lawvere categories \underline{enriched in the category of extended
            metric spaces}.
            \medskip
        \item[Monads on $\mathbf{Set}$ (\underline{EMet})]\hfill\\
            \underline{quantitative} effects.
    \end{description}
    \vfill
    \hfill{\small Mardare, Panangaden and Plotkin (LICS'16, LICS'17);}\\
    \hfill{\small Bacci, Mardare, Panangaden and Plotkin.}
\end{frame}

\begin{frame}{The metric view}
    If $s \qeq{\eps} t$, then for an interpretation
    $\den{\cdot} $ we ask
    \[ d\bigl(\den{s},\den{t}\bigr) \;\leq\; \eps , \]
    where $d$ is a metric.
    \vfill
    \pause
    \begin{columns}[c]
        \begin{column}{0.45\textwidth}
            \[
                \inferrule{a \qeq{\eps} b \\ b \qeq{\delta} c}
                {a \qeq{\eps + \delta} c}
            \]
        \end{column}
        \begin{column}{0.5\textwidth}
            \centering
            \begin{tikzpicture}[scale=1.0]
                \node (a) at (0,0) {$a$};
                \node (c) at (2.2,0) {$c$};
                \node (b) at (1.1,1.1) {$b$};
                \draw (a) -- node[above left,font=\small] {$\eps$} (b);
                \draw (b) -- node[above right,font=\small] {$\delta$} (c);
                \draw (a) -- node[below,font=\small] {$\leq \eps + \delta$} (c);
            \end{tikzpicture}
        \end{column}
    \end{columns}
    \vfill
    Since $d$ is a metric: The triangle inequality makes the rule
    above sound, and $d(a,b) = 0 \Rightarrow a = b$ makes $s \qeq{0} t$
    mean strict equality.
\end{frame}

\begin{frame}{Relaxing to fuzzy relations}
    A generalization relaxes the metric requirement to simply a fuzzy
    relation
    \[ d : A \times A \longrightarrow [0,1] . \]
    \vfill
    No additional assumptions: no symmetry, triangle
    inequality, or $d(a,a) = 0$.
    \vfill
    \pause
    Some caveats with this generalization:
    \begin{itemize}
        \item[--] The "transitivity" rule becomes unsound, and so do
            quantitative reflexivity and symmetry.
        \item[--] We need $=$ as a separate judgement form denoting exact
            equality, since $s \qeq{0} t$ does not imply
            $\den{s} = \den{t}$.
    \end{itemize}
    \vfill
    \hfill{\small Mio, Sarkis and Vignudelli, LMCS 20 (2024).}
\end{frame}

\begin{frame}[fragile]{Formalization: fuzzy spaces and algebras}
    \begin{srcbox}
Record fuzzy_space (R : realType) := {
    fcarrier : Type;
    frel : fcarrier -> fcarrier -> R;
    frel_range : forall x y, (0 <= frel x y <= 1)%R;
}.

Definition frel_nonexpansive {R : realType} (A B : fuzzy_space R)
    (f : fcarrier A -> fcarrier B) : Prop :=
  forall x y, (frel B (f x) (f y) <= frel A x y)%R.

Record FuzzyAlgebra (R : realType) (sig : signature) := {
  fa_space : fuzzy_space R;
  fa_ops : algebra_ops sig (fcarrier fa_space)
}.
    \end{srcbox}
\end{frame}

\begin{frame}{What a quantitative equation looks like}
    Let $\mathbb{A} = (A, \{op^{\mathbb{A}}\}, d_{\mathbb{A}})$ be a
    quantitative algebra.
    \vfill
    \[
        \mathbb{A} \models \forall (B,d_B).\, s = t
        \quad\text{iff}\quad
        \den{s} = \den{t}
    \]
    \vspace{-1em}
    \[
        \mathbb{A} \models \forall (B,d_B).\, s \qeq{\eps} t
        \quad\text{iff}\quad
        d_{\mathbb{A}}\bigl(\den{s},\den{t}\bigr) \leq \eps
    \]
    \vfill
    for all non-expansive $\den{\cdot} : (B,d_B) \to
    (A,d_\mathbb{A})$.
    \vfill
    \pause
    {\small The context $(B,d_B)$ is itself a fuzzy relation on the
    variables. It restricts which assignments are admissible. }
\end{frame}

\begin{frame}[fragile]{Formalization: judgements and interpretations}
    \begin{columns}[T]
        \begin{column}{0.5\textwidth}
            \begin{srcbox}
Inductive judgement {R : realType}
    (sig : signature) (X : Type) :=
  | EqJ  : term sig X -> term sig X ->
           judgement sig X
  | QEqJ : R -> term sig X -> term sig X ->
           judgement sig X.
            \end{srcbox}
        \end{column}
        \begin{column}{0.5\textwidth}
            \begin{srcbox}
Record interpretation
    {R : realType} {sig}
    (X : fuzzy_space R)
    (A : FuzzyAlgebra R sig) := {
  interpretation_fun :>
    fcarrier X -> fcarrier (fa_space A);
  interpretation_nexp :
    frel_nonexpansive X (fa_space A)
      interpretation_fun;
}.
            \end{srcbox}
        \end{column}
    \end{columns}
    \vfill
    \begin{srcbox}
Fixpoint fuzzy_eval {R : realType} {sig X} (A : FuzzyAlgebra R sig)
    (rho : X -> fcarrier (fa_space A)) (t : term sig X) :
    fcarrier (fa_space A) :=
  match t with
  | Var x => rho x
  | App f args => fa_ops A f (fun i => fuzzy_eval A rho (args i))
  end.
    \end{srcbox}
    The two judgement forms of the previous slide are the two
    constructors.
\end{frame}

\begin{frame}[fragile]{Formalization: satisfaction, theories, models}
    \begin{srcbox}
Definition satisfies {R : realType} {sig} (A : FuzzyAlgebra R sig)
    (X : fuzzy_space R) (phi : judgement sig (fcarrier X)) : Prop :=
  forall rho : interpretation X A,
    match phi with
    | EqJ s t => fuzzy_eval A rho s = fuzzy_eval A rho t
    | QEqJ eps s t =>
      (frel (fa_space A) (fuzzy_eval A rho s) (fuzzy_eval A rho t) <= eps)%R
    end.

Definition fuzzy_theory (R : realType) (sig : signature) :=
  forall X : fuzzy_space R, @judgement R sig (fcarrier X) -> Prop.

Definition fuzzy_models {R : realType} {sig} (A : FuzzyAlgebra R sig)
    (E : fuzzy_theory R sig) : Prop :=
  forall X phi, E X phi -> satisfies A X phi.
    \end{srcbox}
    \vfill
    \rid{satisfies} is $\models$ for a single judgement and
    \rid{fuzzy\_models} is $\mathbb{A} \models E$.
\end{frame}

% ============================================================
\section{The quantitative theory $\ICA$}
% ============================================================

\begin{frame}{A quantitative theory of barycentric algebras}
    \begin{align*}
        \text{(Idempotency)}\quad &
        \forall\ctx{c}{ & x \\ x & 1 }.\;
        x \bpl{p} x = x
        \\[0.4em]
        \text{(Skew comm.)}\quad &
        \forall\ctx{cc}{ & x & y \\ x & 1 & 1 \\ y & 1 & 1 }.\;
        x \bpl{p} y = y \bpl{1-p} x
        \\[0.4em]
        \text{(Skew assoc.)}\quad &
        \forall\ctx{ccc}{ & x & y & z \\
        x & 1 & 1 & 1 \\ y & 1 & 1 & 1 \\ z & 1 & 1 & 1 }.\;
        (x \bpl{p} y) \bpl{q} z = x \bpl{pq}
        \bigl(y \bpl{\frac{(1-p)q}{1-pq}} z\bigr)
        \\[0.6em]
        \text{(Interpolative)}\quad &
        \forall\ctx{cccc}{ & x & y & w & z \\
            x & 1 & 1 & \eps & 1 \\
            y & 1 & 1 & 1 & \delta \\
            w & 1 & 1 & 1 & 1 \\
        z & 1 & 1 & 1 & 1 }.\;
        x \bpl{p} y \;\qeq{p\eps + (1-p)\delta}\; w \bpl{p} z
    \end{align*}
    \vfill
    {\footnotesize The first three are the old axioms over the discrete
        context. The fourth constrains only $d(\den{x},\den{w}) \leq \eps$ and
    $d(\den{y},\den{z}) \leq \delta$; every other entry is $1$.}
    \hfill{\footnotesize (the theory $\ICA$)}
\end{frame}

\begin{frame}[fragile]{Formalization: the $\ICA$ signature and its contexts}
    \begin{srcbox}
Record ica_weight (R : realType) := {
  ica_weight_val : R;
  ica_weight_open : (0 < ica_weight_val < 1)%R;   (* strictly inside (0,1) *)
}.
Inductive ica_sym (R : realType) : Type :=
  | ica_plus : ica_weight R -> ica_sym R.
Definition ica_signature {R : realType} : signature :=
  {| sym := ica_sym R; arity := fun _ => 2 |}.
Notation "x <+ p +> y" := (ica_op p x y).
    \end{srcbox}

    \begin{srcbox}
Definition ica_full_space (R : realType) (n : nat) : fuzzy_space R.
  (* 'I_n with frel = fun _ _ => 1 : the discrete fuzzy relation *)

Definition ica_interp_rel {R : realType} (eps delta : R) (a b : 'I_4) : R :=
  match nat_of_ord a, nat_of_ord b with
  | 0, 2 => eps  | 1, 3 => delta  | _, _ => 1
  end.
    \end{srcbox}
    Under $0 = x$, $1 = y$, $2 = w$, $3 = z$ the second one is the
    interpolative matrix, verbatim.
\end{frame}

\begin{frame}[fragile]{Formalization: \rid{Inductive ica\_theory}}
    \begin{srcbox}
Inductive ica_theory (R : realType) : fuzzy_theory R (@ica_signature R) :=
  | ICA_Idem (p : ica_weight R) :
      @ica_theory R (ica_full_space R 1)
        (EqJ ((Var (inord 0)) <+ p +> (Var (inord 0))) (Var (inord 0)))
  | ICA_Skew_Comm (p : ica_weight R) :
      @ica_theory R (ica_full_space R 2)
        (EqJ ((Var (inord 0)) <+ p +> (Var (inord 1)))
             ((Var (inord 1)) <+ ica_weight_complement p +> (Var (inord 0))))
  | ICA_Skew_Assoc (p q : ica_weight R) :
      @ica_theory R (ica_full_space R 3)
        (EqJ (((Var (inord 0)) <+ p +> (Var (inord 1))) <+ q +> (Var (inord 2)))
             ((Var (inord 0)) <+ ica_weight_product p q +>
               ((Var (inord 1)) <+ ica_weight_assoc_inner p q +> (Var (inord 2)))))
  | ICA_Interp (p : ica_weight R) (eps delta : R)
      (Heps : (0 <= eps <= 1)%R) (Hdelta : (0 <= delta <= 1)%R) :
      @ica_theory R (@ica_interp_space R eps delta Heps Hdelta)
        (QEqJ (weighted_sum (ica_probability_weight p) eps delta)
          ((Var (inord 0)) <+ p +> (Var (inord 1)))
          ((Var (inord 2)) <+ p +> (Var (inord 3)))).
    \end{srcbox}
        \vfill
        An \rid{Inductive} in \rid{fuzzy\_theory R sig}, one constructor per
        axiom scheme.
    \end{frame}

    % ============================================================
    \section{The Kantorovich model}
    % ============================================================

    \begin{frame}{An aside about moving piles of sand}
        \begin{center}
            \begin{tikzpicture}[scale=0.8]
                % left pile: two bumps
                \draw[thick,rounded corners=1pt]
                (0,0) .. controls (0.4,0.1) and (0.5,1.4) .. (0.9,1.4)
                .. controls (1.2,1.4) and (1.2,0.5) .. (1.45,0.5)
                .. controls (1.7,0.5) and (1.75,1.5) .. (2.1,1.5)
                .. controls (2.5,1.5) and (2.6,0.1) .. (3.1,0);
                % right pile: one bump
                \draw[thick,rounded corners=1pt]
                (6.2,0) .. controls (6.9,0.1) and (7.1,1.6) .. (7.8,1.6)
                .. controls (8.5,1.6) and (8.7,0.1) .. (9.4,0);
                % figure with shovel
                \begin{scope}[shift={(3.9,0.0)},scale=0.9]
                    \draw[thick] (0.5,1.9) circle (0.28);
                    \fill (0.42,1.95) circle (0.03); \fill
                    (0.58,1.95) circle (0.03);
                    \draw[thick] (0.5,1.62) -- (0.5,0.85);
                    \draw[thick] (0.5,1.45) -- (0.05,1.05);
                    \draw[thick] (0.5,1.45) -- (0.95,1.15);
                    \draw[thick] (0.5,0.85) -- (0.2,0.05);
                    \draw[thick] (0.5,0.85) -- (0.8,0.05);
                    \draw[thick] (0.05,1.05) -- (-0.25,0.35);
                    \draw[thick,fill=goalbg] (-0.42,0.35) -- (-0.08,0.35)
                    -- (-0.13,0.05) -- (-0.37,0.05) -- cycle;
                \end{scope}
                \draw[-{Stealth[length=3mm]},thick] (0.3,-0.45) -- (9.1,-0.45);
                \node at (1.55,-0.95) {$\mu$};
                \node at (7.8,-0.95) {$\nu$};
            \end{tikzpicture}
        \end{center}
        \vfill
        Out of all the possible ways to move the sand, which is the least
        amount of work? (Also known as optimal transport)
        \[
            \Kan(d)(\mu,\nu) \;=\; \inf_{\gamma \in \Cpl(\mu,\nu)}
            \underbrace{\sum_{a,b \in A} \gamma(a,b) \cdot d(a,b)}_{V_d(\gamma)}
        \]
        \vfill
        $\Cpl(\mu,\nu)$ is the set of joint distributions with
        marginals $\mu$ and $\nu$, the transport plans.
    \end{frame}

    \begin{frame}[fragile]{Formalization: couplings and the lifting}
        \begin{srcbox}
Record coupling {R : realType} {X Y : finType}
    (mu : probability_distribution R X)
    (nu : probability_distribution R Y) := {
  coupling_distribution : joint_probability_distribution R X Y;
  coupling_fst : forall x, \sum_(y : Y) coupling_distribution x y = mu x;
  coupling_snd : forall y, \sum_(x : X) coupling_distribution x y = nu y;
}.

Definition coupling_cost {R : realType} {X Y : finType} (d : X -> Y -> R)
    {mu} {nu} (gamma : coupling mu nu) : R :=
  \sum_(x : X) \sum_(y : Y) coupling_distribution gamma x y * d x y.

Definition kantorovich_lifting {R : realType} {X : finType} (d : X -> X -> R)
    (mu nu : probability_distribution R X) : R :=
  inf [set r : R | exists gamma : coupling mu nu, r = coupling_cost d gamma].
        \end{srcbox}
        \begin{itemize}
            \item A \rid{coupling} is a joint distribution bundled with its two
                marginal equations.
        \end{itemize}
    \end{frame}

    \begin{frame}{A model in the quantitative setting}
        The finitely supported distributions, their convex sums, and the
        Kantorovich lifting form a quantitative algebra
        $(\Dist(A),\{\bpl{p}\},\Kan(d))$.
        \vfill
        \pause
        Adding the interpolative axiom is exactly what axiomatises it:
        \vfill
        \begin{block}{Soundness \hfill{\normalfont\small (Prop.~4.6
            and Cor.~4.7)}}
            If $s \qeq{\eps} t$ is derivable, then
            $\Kan(d)(\den{s},\den{t}) \leq \eps$.
        \end{block}
        \vfill
        \begin{block}{Completeness \hfill{\normalfont\small (Prop.~4.10)}}
            If $\Kan(d)(\den{s},\den{t}) \leq \eps$, then $s \qeq{\eps} t$ is
            derivable.
        \end{block}
        \vfill
        \hfill{\small Mio, Compact Quantitative Theories of Convex Algebras,
        MFPS 2025.}
    \end{frame}

    \begin{frame}[fragile]{Formalization: the Kantorovich algebra}
        \begin{srcbox}
Section KantorovichModel.
Context {R : realType} {X : finType}.
Variable d : X -> X -> R.
Hypothesis Hd : forall x y, (0 <= d x y <= 1)%R.

Definition finite_fuzzy_space : fuzzy_space R :=
  {| fcarrier := X; frel := d; frel_range := Hd |}.

Definition kantorovich_fuzzy_space : fuzzy_space R :=
  {| fcarrier := probability_distribution R X;
     frel := kantorovich_lifting d;
     frel_range := fun mu nu => fuzzy_kantorovich_lifting mu nu Hd |}.

Definition kantorovich_algebra : FuzzyAlgebra R (@ica_signature R) :=
  {| fa_space := kantorovich_fuzzy_space; fa_ops := kantorovich_ops |}.
        \end{srcbox}
        \vfill
        \rid{kantorovich\_ops} interprets \rid{ica\_plus p} as
        \rid{convex\_mixture}, and \rid{fuzzy\_kantorovich\_lifting} is the
        bound $0 \leq \Kan(d) \leq 1$, which is the range obligation that
        makes $\Kan(d)$ a legal \rid{fuzzy\_space}. So this record is
        $(\Dist(A),\{\bpl{p}\},\Kan(d))$, and the goal is now
        \rid{fuzzy\_models kantorovich\_algebra ica\_theory}.
    \end{frame}

    \begin{frame}{The proof idea: run both transport plans in parallel}
        Assume $\Kan(d)(\mu_x,\mu_w) \leq \eps$ and
        $\Kan(d)(\mu_y,\mu_z) \leq \delta$. We must show
        \[
            \Kan(d)(\mu_x \bpl{p} \mu_y,\ \mu_w \bpl{p} \mu_z)
            \;\leq\; p\,\eps + (1-p)\,\delta .
        \]
        \pause
        \vfill
        \begin{enumerate}
            \item $\Kan$ is an infimum, so take almost optimal plans:
                $\gamma \in \Cpl(\mu_x,\mu_w)$ and
                $\gamma' \in \Cpl(\mu_y,\mu_z)$ with
                $V_d(\gamma) \leq \eps + \lambda$,
                $V_d(\gamma') \leq \delta + \lambda$.
                \medskip
            \item $\gamma \bpl{p} \gamma' \in
                \Cpl(\mu_x \bpl{p} \mu_y,\ \mu_w \bpl{p} \mu_z)$: run the first
                plan on the $p$-fraction of the mass and the second on the rest.
                \medskip
            \item $V_d$ is "mixture-preserving":
                $V_d(\gamma \bpl{p} \gamma') = p\,V_d(\gamma) +
                (1-p)V_d(\gamma')$.
                \medskip
            \item The infimum is below this one plan, so the bound is
                $p\eps + (1-p)\delta + \lambda$; then let $\lambda \to 0$.
        \end{enumerate}
        \vfill
    \end{frame}

    \begin{frame}[fragile]{The lemmas the proof needs}
        \begin{srcbox}
Lemma coupling_convex_mixture (p : probability_weight R)
    (gamma1 : coupling mu1 nu1) (gamma2 : coupling mu2 nu2) :
  coupling (convex_mixture p mu1 mu2) (convex_mixture p nu1 nu2).
        \end{srcbox}
        \begin{srcbox}
Lemma kantorovich_almost_optimal (d : X -> X -> R) (mu nu : ...) (eta : R) :
  (forall x y, (0 <= d x y)%R) -> (0 < eta)%R ->
  exists gamma : coupling mu nu,
    (coupling_cost d gamma < kantorovich_lifting d mu nu + eta)%R.

Lemma coupling_cost_convex_mixture (p : probability_weight R) ... :
  coupling_cost d (coupling_convex_mixture p gamma1 gamma2)
    = weighted_sum p (coupling_cost d gamma1) (coupling_cost d gamma2).

Lemma kantorovich_convex_mixture (d : X -> X -> R) (p : probability_weight R) :
  (forall x y, (0 <= d x y <= 1)%R) ->
  (kantorovich_lifting d (convex_mixture p mu1 mu2) (convex_mixture p nu1 nu2)
     <= weighted_sum p (kantorovich_lifting d mu1 nu1)
                       (kantorovich_lifting d mu2 nu2))%R.
        \end{srcbox}
            \end{frame}

            \newcommand{\modpane}[1]{%
                \begin{scriptbox}[fontupper=\ttfamily\fontsize{5.8pt}{6.6pt}\selectfont\frenchspacing]%
                    \pstep{#1}\modelscript
            \end{scriptbox}}

            \begin{frame}[fragile]{Modelhood: the three equational goals}
                \begin{columns}[T]
                    \begin{column}{0.44\textwidth}\modpane{2}
                    \end{column}
                    \begin{column}{0.54\textwidth}
                        \begin{goalbox}[title={Goal 1 / 4 \quad (ICA\_Idem)}]
                            \begin{gctx}
d : X -> X -> R
Hd : forall x y : X, 0 <= d x y <= 1
p : ica_weight R
rho : interpretation (ica_full_space R 1)
        kantorovich_algebra
                            \end{gctx}
                            \tcblower
                            \begin{gconc}
|- fuzzy_eval kantorovich_algebra rho
     (Var (inord 0) <+ p +> Var (inord 0)) =
   fuzzy_eval kantorovich_algebra rho (Var (inord 0))
                            \end{gconc}
                        \end{goalbox}
                        \smallskip
                        {\footnotesize \rid{case: Hphi} on the
                            inductive theory gives one
                            goal per axiom, four in total.
                            \\[0.5em]
                            The first three are equalities of
                            distributions. They are
                            closed by the \rid{convex\_mixture\_*}
                            lemmas from earlier, one
                        tactic each.}
                    \end{column}
                \end{columns}
            \end{frame}

            \begin{frame}[fragile]{Modelhood: the interpolative goal}
                \begin{columns}[T]
                    \begin{column}{0.44\textwidth}\modpane{3}
                    \end{column}
                    \begin{column}{0.54\textwidth}
                        \begin{goalbox}[title={Goal 4 / 4 \quad (ICA\_Interp),
                            as it arrives}]
                            \begin{gctx}
p : ica_weight R
eps, delta : R
Heps : 0 <= eps <= 1
Hdelta : 0 <= delta <= 1
rho : interpretation (ica_interp_space Heps Hdelta)
        kantorovich_algebra
                            \end{gctx}
                            \tcblower
                            \begin{gconc}
|- frel (fa_space kantorovich_algebra)
     (fuzzy_eval kantorovich_algebra rho
        (Var (inord 0) <+ p +> Var (inord 1)))
     (fuzzy_eval kantorovich_algebra rho
        (Var (inord 2) <+ p +> Var (inord 3)))
   <= weighted_sum (ica_probability_weight p) eps delta
                            \end{gconc}
                        \end{goalbox}
                        {\footnotesize The context is the four-point
                            space carrying $(0,2) \mapsto \eps$ and
                            $(1,3) \mapsto \delta$. First: evaluate the
                        two $\bpl{p}$-terms.}
                    \end{column}
                \end{columns}
            \end{frame}

            \begin{frame}[fragile]{Modelhood: convexity of the lifting}
                \begin{columns}[T]
                    \begin{column}{0.44\textwidth}\modpane{4}
                    \end{column}
                    \begin{column}{0.54\textwidth}
                        \begin{goalbox}[title={after
                            \texttt{rewrite !kantorovich\_eval\_op}}]
                            \begin{gconc}
|- frel (fa_space kantorovich_algebra)
     (convex_mixture (ica_probability_weight p)
        (fuzzy_eval kantorovich_algebra rho (Var (inord 0)))
        (fuzzy_eval kantorovich_algebra rho (Var (inord 1))))
     (convex_mixture (ica_probability_weight p)
        (fuzzy_eval kantorovich_algebra rho (Var (inord 2)))
        (fuzzy_eval kantorovich_algebra rho (Var (inord 3))))
   <= weighted_sum (ica_probability_weight p) eps delta
                            \end{gconc}
                        \end{goalbox}
                        {\footnotesize Both sides are now literal convex
                            mixtures, which is exactly the shape of
                        \rid{kantorovich\_convex\_mixture}.}
                    \end{column}
                \end{columns}
            \end{frame}

            \begin{frame}[fragile]{Modelhood: split the weighted sum}
                \begin{columns}[T]
                    \begin{column}{0.44\textwidth}\modpane{5}
                    \end{column}
                    \begin{column}{0.54\textwidth}
                        \begin{goalbox}[title={after
                                    \texttt{apply: (le\_trans
                            (kantorovich\_convex\_mixture ...))}}]
                            \begin{gconc}
|- weighted_sum (ica_probability_weight p)
     (kantorovich_lifting d
        (fuzzy_eval kantorovich_algebra rho (Var (inord 0)))
        (fuzzy_eval kantorovich_algebra rho (Var (inord 2))))
     (kantorovich_lifting d
        (fuzzy_eval kantorovich_algebra rho (Var (inord 1)))
        (fuzzy_eval kantorovich_algebra rho (Var (inord 3))))
   <= weighted_sum (ica_probability_weight p) eps delta
                            \end{gconc}
                        \end{goalbox}
                        {\footnotesize The mixture is gone: both sides are
                            $\oplus_p$, and $\oplus_p$ is monotone in each
                        argument.}
                    \end{column}
                \end{columns}
            \end{frame}

            \begin{frame}[fragile]{Modelhood: non-expansiveness of
                \texttt{rho}}
                \begin{columns}[T]
                    \begin{column}{0.44\textwidth}\modpane{6}
                    \end{column}
                    \begin{column}{0.54\textwidth}
                        \begin{goalbox}[title={after
                                    \texttt{apply: weighted\_sum\_le}
                            --- Goal 1 / 2}]
                            \begin{gconc}
|- kantorovich_lifting d
     (fuzzy_eval kantorovich_algebra rho (Var (inord 0)))
     (fuzzy_eval kantorovich_algebra rho (Var (inord 2)))
   <= eps
                            \end{gconc}
                        \end{goalbox}
                        \begin{goalbox}[title={Goal 2 / 2}]
                            \begin{gconc}
|- kantorovich_lifting d
     (fuzzy_eval kantorovich_algebra rho (Var (inord 1)))
     (fuzzy_eval kantorovich_algebra rho (Var (inord 3)))
   <= delta
                            \end{gconc}
                        \end{goalbox}
                        {\footnotesize One goal per component. Each is
                            $\Kan(d)(\den{x},\den{w}) \leq \eps$, and
                        \rid{rho} is non-expansive by construction.}
                    \end{column}
                \end{columns}
            \end{frame}

            \begin{frame}[fragile]{Modelhood: the context matrix, read off}
                \begin{columns}[T]
                    \begin{column}{0.44\textwidth}\modpane{7}
                    \end{column}
                    \begin{column}{0.54\textwidth}
                        \begin{goalbox}[title={after \texttt{apply: (le\_trans
                            (interpretation\_nexp rho (inord 0) (inord 2)))}}]
                            \begin{gconc}
|- frel (ica_interp_space Heps Hdelta) (inord 0) (inord 2) <= eps
                            \end{gconc}
                        \end{goalbox}
                        \begin{goalbox}[title={after
                            \texttt{rewrite /= /ica\_interp\_rel}}]
                            \begin{gconc}
|- match nat_of_ord (inord 0) with
   | 0 => match nat_of_ord (inord 2) with 2 => eps | _ => 1 end
   | 1 => match nat_of_ord (inord 2) with 3 => delta | _ => 1 end
   | _ => 1
   end <= eps
                            \end{gconc}
                        \end{goalbox}
                        {\footnotesize \rid{!inordK} computes the two indices,
                            the matrix returns its $(0,2)$ entry, and the goal
                        is \rid{eps <= eps}.}
                    \end{column}
                \end{columns}
            \end{frame}

            \begin{frame}[fragile]{Modelhood: the second component}
                \begin{columns}[T]
                    \begin{column}{0.44\textwidth}\modpane{8}
                    \end{column}
                    \begin{column}{0.54\textwidth}
                        \begin{goalbox}[title={the one remaining goal}]
                            \begin{gconc}
|- kantorovich_lifting d
     (fuzzy_eval kantorovich_algebra rho (Var (inord 1)))
     (fuzzy_eval kantorovich_algebra rho (Var (inord 3)))
   <= delta
                            \end{gconc}
                        \end{goalbox}
                        {\footnotesize The same two tactics at $(1,3)$, where
                            the matrix returns $\delta$. \\[0.4em]
                            So the whole interpolative axiom is: convexity of
                            the lifting, monotonicity of $\oplus_p$, and two
                        lookups in the context matrix.}
                    \end{column}
                \end{columns}
            \end{frame}

            % ============================================================
            \section{Soundness}
            % ============================================================

            \begin{frame}{The general rules}
                What survives when $d$ is an arbitrary fuzzy relation:
                \vspace{0.6em}
                \begin{columns}[c]
                    \begin{column}{0.34\textwidth}
                        \[
                            \inferrule{ }{x \qeq{d_B(x,y)} y}
                        \]
                        \centering{\footnotesize read off the context}
                    \end{column}
                    \begin{column}{0.32\textwidth}
                        \[
                            \inferrule{\eps \leq \delta \\ a
                            \qeq{\eps} b}{a \qeq{\delta} b}
                        \]
                        \centering{\footnotesize weakening}
                    \end{column}
                    \begin{column}{0.3\textwidth}
                        \[
                            \inferrule{ }{a \qeq{1} b}
                        \]
                        \centering{\footnotesize $d \leq 1$ always}
                    \end{column}
                \end{columns}
                \vspace{1.2em}
                \[
                    \inferrule{\{\, s \qeq{\delta} t \mid \eps < \delta \,\}}
                    {s \qeq{\eps} t}
                \]
                \begin{itemize}
                    \item[--] The bottom rule has infinite
                        width, so proofs are
                        no longer finite objects.
                    \item[--] Known as the order-completeness,
                        Archimedean, or continuity rule.
                    \item[--] Proven to be not needed
                        for barycentric
                        algebras: $\ICA$ is a compact theory.
                \end{itemize}
            \end{frame}

            \begin{frame}[fragile]{Formalization: one inductive, two modes}
                \begin{srcbox}
Inductive frel_derives {R : realType} {sig : signature}
    (E : fuzzy_theory R sig) : frel_derivation_mode ->
    forall X : fuzzy_space R, @judgement R sig (fcarrier X) -> Prop :=
  | FD_Init : forall mode X phi, E X phi -> frel_derives E mode phi
  | FD_EqRefl | FD_EqSym | FD_EqTrans | FD_EqCong  (* Birkhoff, for [=] only *)
  | FD_UseVariables : forall mode X (x y : fcarrier X),
      frel_derives E mode (QEqJ (frel X x y) (Var x) (Var y))
  | FD_Subst : forall mode X Y phi (sigma : fcarrier X -> term sig (fcarrier Y)),
      frel_derives E mode phi ->
      (forall x y, frel_derives E mode (QEqJ (frel X x y) (sigma x) (sigma y))) ->
      frel_derives E mode (subst_judgement sigma phi)
  | FD_QEqReplaceL | FD_QEqReplaceR   (* rewriting a [=] inside a [=_eps] *)
  | FD_Up : forall mode X eps delta s t, (eps <= delta)%R ->
      frel_derives E mode (QEqJ eps s t) -> frel_derives E mode (QEqJ delta s t)
  | FD_Max : forall mode X s t, frel_derives E mode (QEqJ 1 s t)
  | FD_OrderComplete : forall X eps (s t : term sig (fcarrier X)),
      (forall delta, (eps < delta)%R -> frel_derives E FRelFull (QEqJ delta s t)) ->
      frel_derives E FRelFull (QEqJ eps s t).
                \end{srcbox}
                    \begin{itemize}\small
                        \item \rid{FD\_Subst}'s second premise is
                            exactly ``$\sigma$ is
                            provably non-expansive''.
                        \item \rid{FD\_OrderComplete} is only allowed at
                            \rid{FRelFull}, so the
                            finitary system lives inside the same
                            inductive, and compactness
                            is the statement that the two modes agree.
                    \end{itemize}
                \end{frame}

                \newcommand{\sndpane}[1]{%
                    \begin{scriptbox}[fontupper=\ttfamily\fontsize{5.8pt}{6.6pt}\selectfont\frenchspacing]%
                        \pstep{#1}\soundnessscript
                \end{scriptbox}}

                \begin{frame}[fragile]{Soundness of the system:
                    $\vdash\ \Rightarrow\ \models$}
                    \begin{columns}[T]
                        \begin{column}{0.44\textwidth}\sndpane{2}
                        \end{column}
                        \begin{column}{0.54\textwidth}
                            \begin{goalbox}[title={Proof state}]
                                \begin{gctx}
Hmodel : forall X0 phi0, E X0 phi0 ->
  forall rho0 : interpretation X0 A,
    match phi0 with ... end
Hderive : frel_derives E mode phi
rho : interpretation X A
                                \end{gctx}
                                \tcblower
                                \begin{gconc}
|- mode = FRelFinite ->
   match phi with
   | EqJ s t => fuzzy_eval A rho s = fuzzy_eval A rho t
   | QEqJ eps s t =>
     frel (fa_space A) (fuzzy_eval A rho s) (fuzzy_eval A rho t) <= eps
   end
                                \end{gconc}
                            \end{goalbox}
                        \end{column}
                    \end{columns}
                \end{frame}

                \begin{frame}[fragile]{Soundness: the \rid{FD\_UseVariables} case}
                    \begin{columns}[T]
                        \begin{column}{0.42\textwidth}\sndpane{5}
                        \end{column}
                        \begin{column}{0.56\textwidth}
                            {\footnotesize The rule $x \qeq{d_B(x,y)} y$: two
                                variables are as far apart as the context
                            space says.}
                            \begin{goalbox}[title={Goal 6 / 12 \quad
                                (FD\_UseVariables)}]
                                \begin{gctx}
X : fuzzy_space R
x, y : fcarrier X
rho : interpretation X A
                                \end{gctx}
                                \tcblower
                                \begin{gconc}
|- frel (fa_space A) (fuzzy_eval A rho (Var x))
                     (fuzzy_eval A rho (Var y))
   <= frel X x y
                                \end{gconc}
                            \end{goalbox}
                            \begin{goalbox}[title={after \texttt{move => //=}}]
                                \begin{gconc}
|- (frel (fa_space A) (rho x) (rho y) <= frel X x y) = true
                                \end{gconc}
                            \end{goalbox}
                            {\footnotesize \rid{/=} evaluates \rid{fuzzy\_eval}
                                at a \rid{Var}, leaving
                                \rid{frel\_nonexpansive} at $x,y$ --- the
                                record's second field. So
                            \rid{apply: (interpretation\_nexp rho)}.}
                        \end{column}
                    \end{columns}
                \end{frame}

                % \begin{frame}[fragile]{Soundness: the infinitary
                %     rule, in both modes}
                %     \begin{columns}[T]
                %         \begin{column}{0.47\textwidth}
                %             \begin{goalbox}[title={\texttt{derives\_fin\_sound},
                %                 Goal 12/12}]
                %                 \begin{gctx}
                %                     Hmode : FRelFull = FRelFinite
                %                 \end{gctx}
                %                 \tcblower
                %                 \begin{gconc}
                %                     |- frel (fa_space A) (fuzzy_eval A rho s)
                %                     (fuzzy_eval A rho t) <= eps
                %                 \end{gconc}
                %             \end{goalbox}
                %             {\footnotesize \rid{FD\_OrderComplete}
                %                 cannot have been used, since
                %                 it only produces derivations at
                %             \rid{FRelFull}. }
                %         \end{column}
                %         \begin{column}{0.51\textwidth}
                %             \begin{goalbox}[title={\texttt{derives\_full\_sound},
                %                 Goal 12/12}]
                %                 \begin{gctx}
                %                     H0 : forall delta : R, eps < delta ->
                %                     forall rho0 : interpretation X A,
                %                     mode = FRelFull ->
                %                     frel (fa_space A) (fuzzy_eval A rho0 s)
                %                     (fuzzy_eval A rho0 t) <= delta
                %                     Hmode : mode = FRelFull
                %                 \end{gctx}
                %                 \tcblower
                %                 \begin{gconc}
                %                     |- frel (fa_space A) (fuzzy_eval A rho s)
                %                     (fuzzy_eval A rho t) <= eps
                %                 \end{gconc}
                %             \end{goalbox}
                %             \begin{srcbox}
                %                 Lemma le_of_all_strict_upper_bounds
                %                 (a eps : R) :
                %                 (forall delta, (eps < delta)%R -> (a
                %                     % <= delta)%R) ->
                %                     (a <= eps)%R.
                %                     Proof.
                %                     move => H. apply
                %                     /Num.Theory.ler_addgt0Pr => e He.
                %                     apply: H. by rewrite Num.Theory.ltrDl.
                %                     Qed.
                %                 \end{srcbox}
                %             \end{column}
                %         \end{columns}
                %     \end{frame}

                \begin{frame}[fragile]{Soundness for Kantorovich}
                    \begin{srcbox}
Definition dirac_interpretation :
    interpretation finite_fuzzy_space kantorovich_algebra :=
  @Build_interpretation R (@ica_signature R)
    finite_fuzzy_space kantorovich_algebra
    (fun x : X => dirac x)
    (fun x y => kantorovich_dirac_le x y Hd).

Definition ica_term_distribution (s : term (@ica_signature R) X) :=
  fuzzy_eval kantorovich_algebra (fun x : X => dirac x) s.

Theorem finite_kantorovich_soundness
    (s t : term (@ica_signature R) X) (eps : R) :
  @derives_full R (@ica_signature R) (@ica_theory R)
    finite_fuzzy_space (QEqJ eps s t) ->
  (kantorovich_lifting d (ica_term_distribution s)
                         (ica_term_distribution t) <= eps)%R.
Proof.
  move => H.
  apply: (derives_full_sound kantorovich_models_ica H dirac_interpretation).
Qed.
                    \end{srcbox}
                    \vfill
                    Three ingredients, one line: generic
                    soundness, modelhood, and the
                    Dirac interpretation. Dirac is not a
                    restriction.
                \end{frame}

                % ============================================================
                \section{Where the development stands}
                % ============================================================

                \begin{frame}[fragile]{Completeness: what is in place}
                    \textbf{Done:} \file{Signature}, \file{FuzzyRelation},
                    \file{FRelDeduction} (both modes sound),
                    \file{ProbabilityDistribution},
                    \file{KantorovichProperties},
                    \file{ICA}.
                    \medskip

                    \textbf{In progress},
                    \file{BarycentricNormalization}:
                    \begin{srcbox}
Fixpoint term_distribution (s : term (@ica_signature R) X)
  : probability_distribution R X.       (* the denotation [s] *)

Lemma term_distribution_support (s : term (@ica_signature R) X) :
  distribution_support (term_distribution s) = term_vars s.

Lemma condition_recover (mu : probability_distribution R X) (x : X) :
  (mu x < 1) -> x \in distribution_support mu ->
  convex_mixture (ica_probability_weight (ica_weight_clamp (mu x)))
    (dirac x) (condition mu x) = mu.
                    \end{srcbox}
                \end{frame}

                \begin{frame}{How to proceed}
                    \begin{enumerate}
                        \item Normal forms. Iterate
                            \rid{condition\_recover} to
                            write every term as an $n$-ary convex
                            combination
                            $\sum_i p_i x_i$ over its support.
                            \medskip
                        \item For $e(\bm x) =
                            \sum_i p_i x_i$ and
                            tuples $\bm a, \bm b$:
                            \[
                                \ICA \vdash \forall(A,d).\, e(\bm
                                a) \qeq{\eps} e(\bm b),
                                \qquad \eps = \textstyle\sum_i
                                p_i\, d(a_i,b_i),
                            \]
                            by induction on $n$, one instance of
                            the interpolative axiom per
                            step.
                            \medskip
                        \item A
                            coupling $\gamma \in
                            \Cpl(\den{s},\den{t})$ is a recipe
                            for rewriting $s$ and $t$ into
                            the same shape $e(\bm x)$; then (2) gives
                            $s \qeq{V_d(\gamma)} t$, and
                            \rid{FD\_Up} weakens to $\eps$.
                            \medskip
                            % \item \textbf{Finiteness.} Step 3 needs
                            %     one coupling, not the
                            %     infimum: with an optimal coupling the
                            %     proof is finite, giving
                            %     compactness of $\ICA$. \\
                            %     {\small $\hookrightarrow$ requires
                            %         compactness of
                            %         $\Cpl(\mu,\nu)$, the one
                            %         genuinely topological ingredient still
                            %     missing.}
                    \end{enumerate}
                \end{frame}

                % ============================================================
                \section{Outlook: the quantum side}
                % ============================================================

                \begin{frame}{Extending to the quantum setting}
                    In my PhD I want to apply this line of
                    thinking to quantum
                    computing: start with a theory for mixed
                    quantum states, later a
                    theory for quantum programs.
                    \vfill
                    Why mixed states?
                    \begin{itemize}
                        \item[--] They can be viewed as
                            probability distributions over pure
                            states.
                            \medskip
                        \item[--] They are a model of the
                            non-quantitative barycentric
                            axioms: $\rho \bpl{p} \sigma = p\rho
                            + (1-p)\sigma$ is again a
                            density matrix.
                            \medskip
                        \item[--] There is already a notion of
                            quantum optimal transport:
                            \[
                                T_C(\rho,\sigma) \;=\;
                                \inf_{\Gamma \in
                                \mathcal{C}(\rho,\sigma)}
                                \operatorname{tr}(C\Gamma)
                            \]
                            where $\mathcal{C}(\rho,\sigma)$ are
                            the quantum couplings, the
                            bipartite states with marginals
                            $\rho$ and $\sigma$.
                    \end{itemize}
                \end{frame}

                \begin{frame}{Should be easy, right?}
                    \begin{itemize}
                        \item[--] Turns out that the barycentric
                            structure can only reveal
                            the transportation of mass.
                            \medskip
                        \item[--] Mixed states can be coupled in
                            many different ways that
                            transport the mass identically, but
                            give different
                            $\operatorname{tr}(C\Gamma)$.
                            \medskip
                            {\small So the
                                interpolative axiom, stated
                                purely in terms of $\bpl{p}$,
                            cannot separate them. }
                            \medskip
                        \item[--] Currently using a mix of graphs
                            and linear algebra, which
                            seems promising.
                    \end{itemize}
                \end{frame}

                % ============================================================
                % ============================================================

                \begin{frame}{Conclusion}
                    \begin{itemize}
                        \item[--] Quantitative equations $s
                            \qeq{\eps} t$ replace exact
                            equality by bounded distance. Over
                            fuzzy relations the two
                            judgement forms $=$ and $\qeq{\eps}$
                            have to be kept apart, and
                            the triangle rule has to go.
                            \medskip
                        \item[--] Formalized in Rocq/MathComp:
                            signatures and terms, fuzzy
                            algebras, interpretations, and the
                            deduction system in two modes,
                            finitary and order-complete, inside a
                            single inductive.
                            \medskip
                        \item[--] Soundness is proved for both modes
                            (\rid{derives\_fin\_sound},
                            \rid{derives\_full\_sound}), and the
                            Kantorovich lifting on $\Dist(X)$ is
                            proved to be a model of
                            $\ICA$
                            (\rid{kantorovich\_models\_ica}),
                            which gives soundness of
                            $\ICA$ with respect to the Kantorovich distance.
                            \medskip
                        \item[--] Completeness is under way: the
                            normalization machinery
                            (\rid{condition\_recover},
                            \rid{term\_distribution\_support}) is
                            in place; the $n$-ary mixture lemma and
                            the coupling argument are next.
                            \medskip
                        \item[--] Goal: the same story for mixed
                            quantum states, where the
                            barycentric structure alone is not enough.
                    \end{itemize}
                    \vfill
                    \begin{center}\large Thank you. Questions?
                    \end{center}
                \end{frame}

                \appendix

                \begin{frame}{References}
                    \small
                    \begin{itemize}
                        \item R.~Mardare, P.~Panangaden, G.~D.~Plotkin.
                            Quantitative algebraic reasoning. LICS'16.
                        \item R.~Mardare, P.~Panangaden, G.~D.~Plotkin.
                            On the axiomatizability of
                            quantitative algebras. LICS'17.
                        \item M.~Mio, R.~Sarkis, V.~Vignudelli.
                            Universal quantitative algebra for
                            fuzzy relations and
                            generalised metric spaces. LMCS 20(4:19), 2024.
                        \item M.~Mio. Compact quantitative
                            theories of convex algebras.
                            ENTICS, MFPS 2025.
                        \item C.~Villani. Optimal Transport: Old
                            and New. Springer, 2009.
                        \item G.~Bacci, R.~Mardare,
                            P.~Panangaden, G.~D.~Plotkin.
                            Sum and tensor of quantitative
                            effects. LMCS 20(4:9), 2024.
                    \end{itemize}
                \end{frame}

                \end{document}
